\section*{الفصل \ref{ch:b2:orbital-forcing} --- التأثير المداري والعصور الجليدية}

\begin{solution}{exo:b2:orbital-forcing:1}
(أ) $1.0167/0.9833 = 1.0340$: فالأوج أبعد من الحضيض بنسبة \qty{3.4}{\percent}.
(ب) $a(1 + e) - a(1 - e) = 2ea = 0.0334\,a$، أي \qty{3.34}{\percent} من $a$، وهو قريب من نسبة 3.4 في المئة
المذكورة عادةً. (ج) يتناقص ضوء الشمس كمقلوب مربّع المسافة: $1.0340^2 = 1.069$،
أي ضوء شمس أكثر بنحو \qty{7}{\percent} عند الحضيض.
\end{solution}

\begin{solution}{exo:b2:orbital-forcing:2}
(أ) اليوم $0.017236 \times \sin\ang{101.37} = 0.017236 \times 0.9804 = 0.01690$؛ وقبل \num{10000} سنة
$0.019971 \times \sin\ang{293.86} = 0.019971 \times (-0.9145) = -0.01826$؛ وقبل \num{6000} سنة
$0.019249 \times \sin\ang{359.99} \approx 0$، وهي قيم الجداول. (ب) طول الشمس عند
الحضيض هو $\varpi + \ang{180}$: \ang{281.37} اليوم، و\ang{473.86} $-$ \ang{360} $=$ \ang{113.86}
قبل \num{10000} سنة، و\ang{179.99} قبل \num{6000} سنة. (ج) قبل \num{10000} سنة وقع الحضيض
بعد انقلاب يونيو بمقدار \ang{24}، وهو أقرب الثلاثة إليه، في الصيف الشمالي؛
وقبل \num{6000} سنة وقع عند اعتدال سبتمبر، حيث تنعدم \omterm{def:b2:orbital-forcing:elements}{المبادرة المناخية}.
\end{solution}

\begin{solution}{exo:b2:orbital-forcing:3}
(أ) التردّدات تُجمع: $1/25.77 + 1/112 = 0.03880 + 0.00893 = \qty{0.04773}{kyr^{-1}}$، أي دور
قدره $1/0.04773 = 21.0$ ألف سنة. (ب) تقع النتيجة داخل مدى الأدوار الرئيسية،
من 19.0 إلى 23.7 ألف سنة. ولا يستطيع رقم واحد أن يعيد إنتاجها: فدوران المدار
هو نفسه مجموع عدّة حركات بتردّدات مختلفة، تتّحد كلّ منها مع
مبادرة المحور لتعطي دورًا واحدًا من أدوار \omterm{def:b2:orbital-forcing:elements}{المبادرة المناخية}؛ والمتوسط هو
\num{23000} سنة المذكورة عادةً.
\end{solution}

\begin{solution}{exo:b2:orbital-forcing:4}
$\sqrt{1 - 0.0167^2} = 0.99986$. للقيمة \ang{22.1}: $1360 \times 0.37622/(\pi \times 0.99986) =
\qty{162.9}{W.m^{-2}}$؛ وللقيمة \ang{24.5}: $1360 \times 0.41469/(\pi \times 0.99986) =
\qty{179.5}{W.m^{-2}}$. والنسبة $\sin\ang{24.5}/\sin\ang{22.1} = 1.102$: فالمتوسط السنوي عند
القطب يزداد بنسبة \qty{10.2}{\percent}، أي بنحو \qty{16.7}{W.m^{-2}}، بين القيمتين القصويين.
\end{solution}

\begin{solution}{exo:b2:orbital-forcing:5}
(أ) عند انقلاب يونيو $(a/r)^2 = \bigl((1 - e\sin\varpi)/(1 - e^2)\bigr)^2$. اليوم $e\sin\varpi
= 0.016898$ و$(0.983102/0.999703)^2 = 0.9671$، مقابل $1 - 0.033796 = 0.9662$؛
وقبل \num{116000} سنة $e\sin\varpi = 0.043939$ و$(0.956061/0.998065)^2 = 0.9176$، مقابل $1 -
0.087878 = 0.9121$. (ب) عند انقلاب ديسمبر $(a/r)^2 = \bigl((1 + e\sin\varpi)/(1 -
e^2)\bigr)^2$: اليوم $1.0347$ مقابل $1.0338$، وقبل \num{116000} سنة $1.0940$ مقابل $1.0879$.
(ج) الحدود المهملة من رتبة $e^2$: $(1 - x)^2/(1 - e^2)^2 \approx 1 - 2x + x^2 + 2e^2$
مع $x = e\sin\varpi$، أي نحو $3e^2$ هنا. وهي $3 \times 0.0172^2 = 0.0009$ اليوم و$3
\times 0.0440^2 = 0.0058$ قبل \num{116000} سنة، قريبة من الفروق التي وجدناها، $0.0009$
و$0.0055$: فالتقريب أسوأ بست مرّات حين يكون الاختلاف المركزي أكبر بمرّتين ونصف
تقريبًا.
\end{solution}

\begin{solution}{exo:b2:orbital-forcing:6}
عند انقلاب يونيو $\delta = \varepsilon = \ang{23.446}$. $\cos h_0 = -\tan\ang{65}\tan\ang{23.446}
= -2.14451 \times 0.43369 = -0.93006$، فيكون $h_0 = \ang{158.44} = \qty{2.7654}{rad}$. ويدوم النهار
$2 \times 158.44/15 = 21.13$~ساعة (تدور الأرض \ang{15} في الساعة). ومع $\sin\ang{65} = 0.90631$،
و$\cos\ang{65} = 0.42262$، و$\sin\delta = 0.39788$، و$\cos\delta = 0.91744$ و$\sin h_0 = 0.36742$،
يكون القوس $2.7654 \times 0.90631 \times 0.39788 + 0.42262 \times 0.91744 \times 0.36742 = 0.99720
+ 0.14246 = 1.13966$. وعند المسافة المتوسطة $\bar Q = (1360/\pi) \times 1.13966 = 432.90 \times 1.13966
= \qty{493.4}{W.m^{-2}}$؛ ومع عامل المسافة في \cref{exo:b2:orbital-forcing:5}، $493.4 \times
0.9671 = \qty{477.1}{W.m^{-2}}$. وللتحقّق، تعطي الخطوات نفسها عند \ang{60} شمالًا \qty{475.96}{W.m^{-2}}، حيث تعطي الجداول \num{475.95}.
\end{solution}

\begin{solution}{exo:b2:orbital-forcing:7}
(أ) اعتدال مارس هو $\lambda = 0$، واعتدال سبتمبر $\lambda = \pi$. ومع $\sqrt{1 - e^2}
\approx 1$، تكون المدّة $\frac{T}{2\pi}\int_0^{\pi}\bigl(1 + 2e\cos(\lambda -
\varpi)\bigr)\,\mathrm d\lambda = \frac{T}{2\pi}\bigl(\pi + 2e[\sin(\lambda - \varpi)]_0^{\pi}\bigr)
= \frac{T}{2\pi}(\pi + 4e\sin\varpi) = \frac T2 + \frac{2T}{\pi}e\sin\varpi$. (ب) نصف السنة الآخر
أقصر بالمقدار نفسه، فيكون الفرق $(4T/\pi)\,e\sin\varpi = (4 \times
365.2422/\pi) \times 0.016898 = 465.04 \times 0.016898 = 7.86$~يوم. (ج) الربيع والصيف معًا
يزيدان على الخريف والشتاء بمقدار $4.5 + 3 = 7.5$~يوم؛ والفرقان المذكوران مقرّبان،
والتوافق ضمن حدود هذا التقريب.
\end{solution}

\begin{solution}{exo:b2:orbital-forcing:8}
(أ) للقيمة $e = 0.0046$: $\sqrt{1 - e^2} = 0.999989$ و$340/0.999989 = \qty{340.004}{W.m^{-2}}$. وللقيمة
$e = 0.0565$: $\sqrt{1 - 0.003192} = 0.998403$ و$340/0.998403 = \qty{340.544}{W.m^{-2}}$. (ب)
الفرق \qty{0.54}{W.m^{-2}}، أي \qty{0.16}{\percent}. (ج) تأرجح إشماس منتصف يوليو عند \ang{65} شمالًا
بمقدار $497 - 388 = \qty{109}{W.m^{-2}}$، أي نحو \qty{25}{\percent} من متوسطه ومئتي
ضعف تغيّر المتوسط العالمي السنوي: فالمدار يؤثّر في المناخ بنقل ضوء الشمس بين
الفصول ودوائر العرض، لا بتغيير المجموع.
\end{solution}

\begin{solution}{exo:b2:orbital-forcing:9}
مع $\omega = 2\pi/P$ و$\tau = \qty{10}{kyr}$: للدور $P = 23$، $\omega\tau = 2.732$، أي سعة
نسبية $1/\sqrt{1 + 7.464} = 0.344$ وتأخّر $\arctan(2.732)/\omega = 1.2199/0.2732 = 4.47$~ألف سنة؛
وللدور $P = 41$، $\omega\tau = 1.532$، والسعة $0.546$، والتأخّر $0.9926/0.1532 = 6.48$~ألف سنة؛ وللدور $P =
100$، $\omega\tau = 0.628$، والسعة $0.847$، والتأخّر $0.5610/0.0628 = 8.93$~ألف سنة. فالغطاء الجليدي
يفضّل التأثير البطيء: إذ يتابع \qty{85}{\percent} من دورة ذات \num{100000} سنة، لكنه لا يتابع إلا
ثلث دورة مبادرة، لأنه لا يملك وقتًا كافيًا لينمو أو ينكمش قبل أن ينعكس
التأثير السريع. فهو يعمل عمل مرشّح تمرير منخفض.
\end{solution}

\begin{solution}{exo:b2:orbital-forcing:10}
(أ) المعادلة خطّية: إذا كان $V_1$ استجابة التأثير $A_1\sin\omega_1t$ و$V_2$ استجابة التأثير
$A_2\sin\omega_2t$، فإن $V_1 + V_2$ استجابة لمجموعهما. وبحسب \cref{prop:b2:orbital-forcing:lag} فإن $V_1$
و$V_2$ دالّتان جيبيتان تردّداهما $\omega_1$ و$\omega_2$ (بعد أن يتلاشى الطور العابر)، فلا
تحمل الاستجابة أيّ تردّد آخر. (ب) مع $V \approx B_1\sin(\omega_1t - \phi_1) +
B_2\sin(\omega_2t - \phi_2)$،
\[
\begin{aligned}
  V^2 &= \tfrac12(B_1^2 + B_2^2) - \tfrac12B_1^2\cos(2\omega_1t - 2\phi_1) -
  \tfrac12B_2^2\cos(2\omega_2t - 2\phi_2)\\
  &\quad + B_1B_2\bigl[\cos\bigl((\omega_1 - \omega_2)t - \phi_1 + \phi_2\bigr) -
  \cos\bigl((\omega_1 + \omega_2)t - \phi_1 - \phi_2\bigr)\bigr],
\end{aligned}
\]
فإن الحدّ $kV^2$ يؤثّر في المعادلة عند تردّدَي الفرق والمجموع، وعند التردّدين
المضاعفين. (ج) $1/(1/19.0 - 1/23.7) = 1/(0.05263 - 0.04219) = 95.8$ ألف سنة؛ والزوج 19.0
و22.4 يعطي 125 ألف سنة. فتردّدات الفرق لحدود المبادرة تقع في
نطاق الاختلاف المركزي، كما يجب، لأن \omterm{def:b2:orbital-forcing:elements}{المبادرة المناخية} هي $e\sin\varpi$. فغطاء جليدي غير خطّي
يستطيع إذن أن يستخرج إيقاعًا من \num{100000} سنة من المبادرة؛ أمّا الخطّي فلا
يستطيع. والدورات ذات \num{100000} سنة دليل على سلوك غير خطّي، وإن كان هذا الحساب
وحده لا يبيّن أيّ سلوك هو.
\end{solution}

\begin{solution}{exo:b2:orbital-forcing:11}
(أ) كما في برهان \cref{prop:b2:orbital-forcing:annual}، $(a/r)^2\,\mathrm dt =
T\,\mathrm d\lambda/(2\pi\sqrt{1 - e^2})$، و$\sin\delta = \sin\varepsilon\sin\lambda$ موجبة
عندما $0 < \lambda < \pi$. فيكون $E = \frac{S_0T\sin\varepsilon}{2\pi\sqrt{1 - e^2}}\int_0^{\pi}\sin\lambda
\,\mathrm d\lambda = S_0T\sin\varepsilon/(\pi\sqrt{1 - e^2})$. (ب) $T = 365.25 \times 86400 =
\qty{3.1558e7}{s}$، و$S_0T/\pi = \qty{1.3661e10}{J.m^{-2}}$؛ وللقيمة \ang{22.1}، $E = 1.3661 \times 10^{10}
\times 0.37622/0.99986 = \qty{5.14}{GJ.m^{-2}}$؛ وللقيمة \ang{24.5}، $\qty{5.67}{GJ.m^{-2}}$،
أي أكثر بنسبة \qty{10}{\percent}. (ج) المبادرة تحدّد كيف تتوزّع الطاقة نفسها في الزمن: فالصيف
عند الحضيض أقصر وأشدّ، والصيف عند الأوج أطول وأضعف. والذوبان لا يبدأ
إلا فوق \qty{0}{\celsius}، فيمكن لصيف قصير شديد أن يُذيب أكثر من صيف طويل بارد
يتلقّى المجموع نفسه.
\end{solution}

\begin{solution}{exo:b2:orbital-forcing:12}
(أ) $3.93\log_2(279.4/189.7) = 3.93\log_2 1.4729 = 3.93 \times 0.5586 = \qty{2.20}{W.m^{-2}}$. (ب)
أصغر ارتفاع هو من 194 إلى \qty{266}{ppm}: $3.93\log_2 1.3711 = 3.93 \times 0.4554 =
\qty{1.79}{W.m^{-2}}$؛ وأكبرها من 188 إلى \qty{282}{ppm}: $3.93\log_2 1.5 = 3.93 \times 0.5850 =
\qty{2.30}{W.m^{-2}}$. (ج) تأثير ثنائي أكسيد الكربون أربعة أضعاف أكبر تغيّر في
الإشماس العالمي السنوي، \qty{0.54}{W.m^{-2}}، وهو يؤثّر في كلّ مكان وطوال السنة؛ أمّا الإشماس
الصيفي عند \ang{65} شمالًا فيتأرجح بنحو \qty{109}{W.m^{-2}}، لكن عند عرض واحد وفي فصل واحد.
فالمدار هو الزناد، عبر صيف الأغطية الجليدية الشمالية؛ وثنائي أكسيد الكربون، مع
\omterm{def:b2:precambrian-earth:ice-albedo}{التغذية الراجعة للجليد والألبيدو}، مضخِّم ينشر التغيّر على الكرة الأرضية.
\end{solution}

\begin{solution}{pb:b2:orbital-forcing:1}
\textbf{1.} آنذاك $0.0440 \times \sin\ang{92.7} = 0.0440 \times 0.99889 = 0.04395$؛ والآن $0.0172 \times
\sin\ang{101.4} = 0.0172 \times 0.98027 = 0.01686$، أي أقلّ بنحو 2.6 مرّة.
\textbf{2.} آنذاك $1 - 0.0440 = 0.956$ و$1 + 0.0440 = 1.044$؛ والآن $0.9828$ و$1.0172$.
\textbf{3.} عند $\lambda = \varpi + \ang{180}$: آنذاك \ang{272.7}، والآن \ang{281.4}، وكلاهما قرب
انقلاب ديسمبر ($\lambda = \ang{270}$)؛ وآنذاك مرّت الأرض بالحضيض بعده بأقلّ من ثلاثة أيام،
وبالأوج عند انقلاب يونيو.
\textbf{4.} آنذاك $(1.044/0.956)^2 = 1.0921^2 = 1.193$، والآن $(1.0172/0.9828)^2 = 1.0350^2 = 1.071$:
أي ضوء شمس أكثر بنسبة \qty{19}{\percent} عند الحضيض آنذاك، و\qty{7}{\percent} الآن.
\textbf{5.} أصغر، \ang{22.52} مقابل \ang{23.45}: فقد كانت شمس الصيف أخفض في العروض العليا،
وتلقّت العروض العليا ضوء شمس أقلّ في الصيف وعلى مدى السنة.
\textbf{6.} آنذاك $340/\sqrt{1 - 0.001936} = 340/0.999032 = \qty{340.33}{W.m^{-2}}$؛ والآن $340/0.999852
= \qty{340.05}{W.m^{-2}}$. والفرق، \qty{0.28}{W.m^{-2}} أي أقلّ من جزء في الألف،
غير معتبر.
\textbf{7.} عند $\lambda = \ang{90}$، $\cos(\ang{90} - \varpi) = \sin\varpi$ في
\cref{prop:b2:orbital-forcing:distance}.
\textbf{8.} آنذاك $(1 - 0.001936)/(1 - 0.04395) = 0.998064/0.956049 = 1.04395$؛ والآن $(1 -
0.000296)/(1 - 0.01686) = 0.999704/0.983139 = 1.01685$.
\textbf{9.} آنذاك $1/1.04395^2 = 0.91758$؛ والآن $1/1.01685^2 = 0.96714$.
\textbf{10.} آنذاك $1 - 2 \times 0.04395 = 0.91210$؛ والآن $1 - 2 \times 0.01686 = 0.96628$: فخطأ
التقريب $0.0055$ آنذاك، و$0.0009$ الآن.
\textbf{11.} آنذاك أضعف بمقدار $1 - 0.91758 = \qty{8.2}{\percent}$، والآن بمقدار \qty{3.3}{\percent}.
\textbf{12.} عند الانقلاب $\delta = \varepsilon$: \ang{22.52} آنذاك، و\ang{23.45} الآن.
\textbf{13.} آنذاك $\cos h_0 = -2.14451 \times 0.41462 = -0.88916$، و$h_0 = \ang{152.77} =
\qty{2.6663}{rad}$؛ والآن $\cos h_0 = -2.14451 \times 0.43378 = -0.93023$، و$h_0 = \ang{158.47} =
\qty{2.7659}{rad}$.
\textbf{14.} $2h_0/\ang{15}$ لكلّ ساعة: آنذاك $20.37$~ساعة، والآن $21.13$~ساعة.
\textbf{15.} آنذاك $2.6663 \times 0.90631 \times 0.38301 + 0.42262 \times 0.92375 \times 0.45759 =
0.92554 + 0.17864 = 1.10418$ (\num{1.10417} من دون تقريب)، و$432.90 \times 1.10417 =
\qty{478.0}{W.m^{-2}}$؛ والآن $2.7659 \times 0.90631 \times 0.39795 + 0.42262 \times 0.91741 \times
0.36697 = 0.99757 + 0.14228 = 1.13985$ (\num{1.13982} من دون تقريب)، و$432.90 \times 1.13982 = \qty{493.4}{W.m^{-2}}$.
\textbf{16.} آنذاك $478.0 \times 0.91758 = \qty{438.6}{W.m^{-2}}$؛ والآن $493.4 \times 0.96714 =
\qty{477.2}{W.m^{-2}}$، أي أكثر بمقدار \qty{38.6}{W.m^{-2}} ممّا كان آنذاك.
\textbf{17.} في منتصف يوليو يكون طول الشمس \ang{120}: $\sin\delta = \sin\varepsilon\sin\ang{120}$،
أي ميل يقارب \ang{20} بدلًا من \ang{23}، فيكون النهار أقصر والشمس أخفض: فعند
\ang{65} شمالًا يكون الإشماس اليومي المتوسط في أعظم قيمه قرب الانقلاب، حين يكون الميل في أعظم قيمه.
\textbf{18.} مع مسافة اليوم، يكلّف تغيّر ميل المحور وحده $(493.4 - 478.0) \times
0.96714 = \qty{14.9}{W.m^{-2}}$؛ ومع ميل المحور القديم، تكلّف المسافة وحدها $478.0 \times
(0.96714 - 0.91758) = \qty{23.7}{W.m^{-2}}$. ومجموعهما \qty{38.6}{W.m^{-2}}؛ فقد كان للمسافة
نحو ثلاثة أخماس الوزن.
\textbf{19.} بعد تغيّر مفاجئ يكون الغطاء الجليدي قد أنجز الكسر $1 - \mathrm e^{-t/\tau}$ منه:
$1 - \mathrm e^{-0.5} = 0.39$ بعد \qty{5}{kyr}، و$1 - \mathrm e^{-1} = 0.63$ بعد \qty{10}{kyr}،
و$1 - \mathrm e^{-2} = 0.86$ بعد \qty{20}{kyr}.
\textbf{20.} $\omega\tau = 2\pi \times 10/23 = 2.732$: يتابع الغطاء الجليدي $1/\sqrt{1 + 2.732^2} =
0.34$ من التأثير، متأخّرًا بمقدار $\arctan(2.732)/\omega = 4.5$~ألف سنة.
\textbf{21.} نصف دورة مبادرة، أي نحو 11 إلى \qty{12}{kyr}، يساوي تقريبًا زمن استجابة واحدًا: فقبل
أن يقطع الغطاء الجليدي ثلثي الطريق، ينعكس الإشماس ويشدّه إلى الوراء.
\textbf{22.} في النموذج الخطّي ينمو الجليد وينكمش بالمعدّل نفسه لتأثيرات من الحجم
نفسه: فتأثير صاعد وتأثير هابط يعطيان استجابتين متناظرتين كصورة في مرآة. أمّا تراكم بطيء يليه
انهيار سريع فيحتاج إلى استجابة تتوقّف على الحالة أو على اتجاه التغيّر، أي إلى
استجابة غير خطّية.
\textbf{23.} ارتفع ثنائي أكسيد الكربون بما يعطي تأثيرًا بين $3.93\log_2(266/194) = 1.8$ و$3.93\log_2(282/188) =
\qty{2.3}{W.m^{-2}}$ من \omterm{def:b2:orbital-forcing:glacials}{الفترة الجليدية} إلى \omterm{def:b2:orbital-forcing:glacials}{الفترة بين الجليدية}: وهو تأثير على الكرة الأرضية كلّها وطوال
السنة. فإذ ينطلق حين تضعف الدورة الأطلسية، ينقل تغيّر الصيف
الشمالي إلى الكوكب كلّه، ويضخّمه مع ألبيدو الجليد المنكمش.
\textbf{24.} حين انتهت آخر فترة بين جليدية، كان الإشماس اليومي المتوسط عند انقلاب يونيو عند
\ang{65} شمالًا نحو \textbf{\qty{439}{W.m^{-2}}}، أي أقلّ بمقدار \qty{39}{W.m^{-2}} من اليوم، وهو أدنى
قيمه في الدورة.
\end{solution}
